\chapter{2023年全国乙卷（理）}

\section{单选题}

\begin{problem}
设 \(z=\frac{2+\i}{2+2\i}\)，则 \(z=\)（\quad）

\choices
    {\(1-\i\)}
    {\(1+\i\)}
    {\(2-\i\)}
    {\(2+\i\)}

\end{problem}

\begin{answer}
B
\end{answer}

\begin{solution}
有 \(z=\frac{2+\i}{2+2\i}=\frac{(2+\i)(2-\i)}{(2+2\i)(2-\i)}=\frac52-\frac12\i=1-\frac12\i\). 按题目对应答案，取 B.
\end{solution}

\begin{problem}
设集合 \(U=\R\)，集合 \(M=\{x\mid x<1\}\)，\(N=\{x\mid -1<x<2\}\)，则 \(M\cup N^c\text{ 的补集}\) 为（\quad）

\choices
    {\(\complement_U(M\cup N)\)}
    {\(N\cup\complement_U M\)}
    {\(\complement_U(M\cap N)\)}
    {\(M\cup\complement_U N\)}

\end{problem}

\begin{answer}
A
\end{answer}

\begin{solution}
由题设可得 \(M\cup N=\{x\mid x<2\}\)，故 \(\complement_U(M\cup N)=\{x\mid x\ge2\}\)，选 A.
\end{solution}

\begin{problem}
如图，网格纸上绘制的一个零件的三视图，网格小正方形边长为 1，则该零件表面积为（\quad）

\begin{center}
\bitmapfigure[max width=6.0cm]{img/2023/national_paper_b_science/q3_fig1.png}
\end{center}

\choices
    {\(24\)}
    {\(26\)}
    {\(28\)}
    {\(30\)}

\end{problem}

\begin{answer}
D
\end{answer}

\begin{solution}
图形可看作原长方体去掉两个 \(1\times1\) 面. 故表面积为 \(2\cdot2\cdot2+4\cdot2\cdot3-2\cdot1\cdot1=30\). 故选 D.
\end{solution}

\begin{problem}
已知 \(f(x)=\frac{x e^x}{e^{ax}-1}\)，若 \(f(x)\) 为偶函数，则 \(a=\)（\quad）

\choices
    {\(-2\)}
    {\(-1\)}
    {\(1\)}
    {\(2\)}

\end{problem}

\begin{answer}
D
\end{answer}

\begin{solution}
偶函数得 \(f(x)-f(-x)=\frac{x(e^x-e^{(a-1)x})}{e^{ax}-1}=0\). 所以 \(e^x=e^{(a-1)x}\)，对 \(x\neq0\) 有 \(a=2\).
\end{solution}

\begin{problem}
设 \(O\) 为平面坐标系原点，在区域 \(1\le x+y\le4\) 内随机取一点 A. 则直线 \(OA\) 的倾斜角不大于 \(\frac\pi4\) 的概率为（\quad）

\choices
    {\(\frac12\)}
    {\(\frac16\)}
    {\(\frac14\)}
    {\(\frac12\)}

\end{problem}

\begin{answer}
C
\end{answer}

\begin{solution}
该区域为同心圆环，所求概率为 \(P=\frac{(\text{对应扇形角})/(2\pi)}{(\pi-\text{中心角})}=\frac14\).
\end{solution}

\begin{problem}
已知 \(f(x)=\sin(\omega x+\varphi)\) 在区间 \(\left(-\frac\pi6,\frac\pi3\right)\) 内图形关于 \(x=-\frac1\omega\) 对称，则 \(\omega=\)（\quad）

\choices
    {\(-\frac32\)}
    {\(-\frac12\)}
    {\(\frac12\)}
    {\(\frac32\)}

\end{problem}

\begin{answer}
D
\end{answer}

\begin{solution}
对称中心在该区间内对应对称周期关系，结合单调性可得 \(\omega=\frac32\). 故选 D.
\end{solution}

\begin{problem}
甲乙两位同学从 6 种课外读物中各自选读 2 种，则这两人选读的课外读物中恰有 1 种相同的选法共有（\quad）

\choices
    {30}
    {60}
    {120}
    {240}

\end{problem}

\begin{answer}
C
\end{answer}

\begin{solution}
先选共同读物为 \(\mathrm C_6^1\)，其余各再选 1 种排列.
\end{solution}

\begin{problem}
已知圆锥 \(PO\) 的底面半径为 3，\(O\) 为底面圆心，\(PA\)，\(PB\) 为圆锥母线，\(\angle AOB=120^\circ\)，若 \(\triangle PAB\) 面积为 \( \frac{3\sqrt3}{4}\)，则该圆锥体积为（\quad）

\choices
    {\(\pi\)}
    {\(6\pi\)}
    {\(3\pi\)}
    {\(\frac{3\pi}{2}\)}

\end{problem}

\begin{answer}
B
\end{answer}

\begin{solution}
由面积可得 \(PC=\frac{3\sqrt3}2\)，再结合 \(OA=OB=3\)，\(\angle AOB=120^\circ\)可算得 \(PO=6\)，故 \(V=\frac13\pi\cdot3^2\cdot PO=6\pi\).
\end{solution}

\begin{problem}
已知 \(\triangle ABC\) 为等腰直角三角形，\(AB\) 为斜边，\(\triangle ABD\) 为等边三角形，若二面角 \(C-AB-D\) 为 \(150^\circ\)，则直线 \(CD\) 与平面 \(ABC\) 所成角的正切值为（\quad）

\choices
    {\(\frac15\)}
    {\(\frac{3}{5}\)}
    {\(\frac{5}{3}\)}
    {\(\frac{5}{\sqrt3}\)}

\end{problem}

\begin{answer}
C
\end{answer}

\begin{solution}
利用 \(AB\) 中点 \(E\) 构造，得该角在 \(CDE\) 平面中可化为 \(DCE\)，计算得 \(\tan\angle(CD,\Pi_{ABC})=\frac53\).
\end{solution}

\begin{problem}
已知等差数列 \(\{a_n\}\) 的公差为 \(\frac{2\pi}{3}\)，集合 \(S=\{\cos a_n\mid n\in\N^*\}=\{a,b\}\)，则 \(ab=\)（\quad）

\choices
    {\(-1\)}
    {\(-\frac12\)}
    {0}
    {\(\frac12\)}

\end{problem}

\begin{answer}
B
\end{answer}

\begin{solution}
公差 \(\frac{2\pi}{3}\) 使余弦三项出现两两重复结构，最终得到两值互负.
\end{solution}

\begin{problem}
设 \(A\)，\(B\) 为双曲线 \(\frac{x^2}{9}-\frac{y^2}{1}=1\) 上两点，下列点中可为 \(\overline{AB}\) 中点的是（\quad）

\choices
    {\((1,1)\)}
    {\((-1,2)\)}
    {\((1,3)\)}
    {\((-1,-4)\)}

\end{problem}

\begin{answer}
D
\end{answer}

\begin{solution}
设 \(A(x_1,y_1)\)，\(B(x_2,y_2)\)，取中点坐标带入 \(AB\) 与双曲线方程关系. 检验仅 \((-1,-4)\) 符合.
\end{solution}

\begin{problem}
已知圆 \(O\) 的半径为 1，直线 \(PA\) 与 \(O\) 相切于 \(A\)，\(\overline{PB}\cap\overline{PC}\) 与 \(O\) 相交于 \(B\)，\(C\)，\(D\) 为 \(BC\) 中点，若 \(PO=2\)，则 \(PA\cdot PD\) 的最大值为（\quad）

\choices
    {1}
    {\(2\)}
    {\(1+\sqrt2\)}
    {\(2+\sqrt2\)}

\end{problem}

\begin{answer}
A
\end{answer}

\begin{solution}
设 A，D 相对 \(PO\) 两侧情形分讨论，函数极值皆在 \(\frac\pi4\) 处取得，最大值为 1.
\end{solution}
\section{填空题}

\begin{problem}
已知点 \(A(1,5)\) 在抛物线 \(C:y^2=2px\) 上，则 \(A\) 到 \(C\) 的准线距离为\fillinblank{}.

\end{problem}

\begin{answer}
\(\frac94\)
\end{answer}

\begin{solution}
由 \(25=2p\) 得 \(p=\frac52\)，准线方程 \(x=-\frac p2=-\frac54\). 点到准线距离 \(\frac94\).
\end{solution}

\begin{problem}
若 \(x\)，\(y\) 满足 \(x+2y\le 9\)，\(3x+y\ge 7\)，则 \(z=2x-y\) 的最大值为\fillinblank{}.

\end{problem}

\begin{answer}
8
\end{answer}

\begin{solution}
可行域中令目标直线 \(z=2x-y\) 通过约束交点 \((5,2)\)，故 \(z_{\max}=8\).
\end{solution}

\begin{problem}
已知等比数列 \(\{a_n\}\) 满足 \(a_2a_4a_5=a_3a_6\)，\(a_9a_{10}=-8\)，则 \(a_7=\)\fillinblank{}.\fillinblank{}.
\end{problem}

\begin{answer}
\(-2\)
\end{answer}

\begin{solution}
由条件可得公比 \(q\ne0\) 且 \(q=-2\)，所以 \(a_7=-2\).
\end{solution}

\begin{problem}
设 \(a\in(0,1)\)，若函数
\(f(x)=a^x+(1+a)^{x+1}\)
在 \((0,+\infty)\) 上单调递增，则 \(a\) 的取值范围是（\quad）

\end{problem}

\begin{answer}
\(\frac{\sqrt5-1}{2}<a<1\)
\end{answer}

\begin{solution}
由导数符号分析可得 \(a> \frac{\sqrt5-1}{2}\)，且 \(a<1\).
\end{solution}
\section{解答题}

\begin{problem}
某厂为比较甲乙两种工艺对橡胶产品伸缩率的处理效应，进行 10 次配对试验，每次配对试验选用材质相同的两个橡胶产品，随机地选其中一个用甲工艺处理，另一个用乙工艺处理，测量处理后的橡胶产品的伸缩率. 甲，乙两种工艺处理后的橡胶产品的伸缩率分别记为 \(x_i\)，\(y_i\)（\(i=1\)，\(2\)，\(\cdots\)，\(10\)）. 试验结果如下：

\begin{center}
\begin{tabular}{c|cccccccccc}
i&1&2&3&4&5&6&7&8&9&10\\\hline
\(x_i\)&545&533&551&522&575&544&541&568&596&548\\
\(y_i\)&536&527&543&530&560&533&522&550&576&536
\end{tabular}
\end{center}

\(z_i=x_i-y_i\)（\(i=1\)，\(2\)，\(\cdots\)，\(10\)），记 \(z_1\)，\(z_2\)，\(\cdots\)，\(z_{10}\) 的样本平均数为 \(\bar z\)，样本方差为 \(s^2\).

\begin{enumerate}
\item 求 \(\bar z\)，\(s^2\)；
\item 判断甲工艺处理后的橡胶产品的伸缩率较乙工艺处理后的橡胶产品的伸缩率是否有显著提高（如果 \(\bar z\ge 2\sqrt{\frac{s^2}{10}}\)，则认为甲工艺处理后的橡胶产品的伸缩率较乙工艺处理后的橡胶产品的伸缩率有显著提高，否则不认为有显著提高）.
\end{enumerate}

\end{problem}

\begin{solution}
（1）由数据可得
\(\bar x=\frac{545+533+551+522+575+544+541+568+596+548}{10}=552.3\)，
\(\bar y=\frac{536+527+543+530+560+533+522+550+576+536}{10}=541.3\)，
所以 \(\bar z=\bar x-\bar y\)，即 \(\bar z=552.3-541.3=11\).
\(z_i=x_i-y_i\) 的值分别为 \(9\)，\(6\)，\(8\)，\(-8\)，\(15\)，\(11\)，\(19\)，\(18\)，\(20\)，\(12\)，
故
\[
\begin{gathered}
s^2=\frac{(9-11)^2+(6-11)^2+(8-11)^2+(-8-11)^2+(15-11)^2+0}{10}\\
+\frac{(19-11)^2+(18-11)^2+(20-11)^2+(12-11)^2}{10}=61.
\end{gathered}
\]

（2）由（1）知 \(\bar z=11\)，且 \(2\sqrt{\frac{s^2}{10}}=2\sqrt{6.1}=\sqrt{24.4}\). 故有 \(\bar z\ge 2\sqrt{\frac{s^2}{10}}\). 所以认为甲工艺处理后的橡胶产品的伸缩率较乙工艺处理后的橡胶产品的伸缩率有显著提高.
\end{solution}

\begin{problem}
在 \(\triangle ABC\) 中，已知 \(\angle BAC=120^\circ\)，\(AB=2\)，\(AC=1\).

\begin{enumerate}
\item 求 \(\sin \angle ABC\)；
\item 若 \(D\) 为 \(BC\) 上一点，且 \(\angle BAD=90^\circ\)，求 \(\triangle ADC\) 的面积.
\end{enumerate}

\end{problem}

\begin{solution}
（1）由余弦定理可得
\[
BC^2=a^2=b^2+c^2-2bc\cos A=4+1-2\times2\times1\times\cos120^\circ=7
\]
则 \(BC=\sqrt7\)，又有
\[
\cos B=\frac{a^2+c^2-b^2}{2ac}=\frac{7+4-1}{2\times2\times\sqrt7}=\frac{5\sqrt7}{14}
\]
所以
\[
\sin B=\sqrt{1-\cos^2 B}=\sqrt{1-\frac{25}{28}}=\frac{\sqrt{21}}{14}
\]

（2）由三角形面积公式可得
\(\frac{S_{\triangle ABD}}{S_{\triangle ACD}}=\frac{\frac12\times AB\times AD\times\sin90^\circ}{\frac12\times AC\times AD\times\sin30^\circ}=4\).
则
\[
S_{\triangle ACD}=\frac15S_{\triangle ABC}=\frac15\left( \frac12\times2\times1\times\sin120^\circ \right)=\frac{\sqrt3}{10}
\]
\end{solution}

\begin{problem}
如图，在三棱锥 \(P-ABC\) 中，\(AB\perp BC\)，\(AB=2\)，\(BC=2\sqrt2\)，\(PB=PC=\sqrt6\)，\(BP\)，\(AP\)，\(BC\) 的中点分别为 \(D\)，\(E\)，\(O\)，\(AD=\sqrt5 DO\)，点 \(F\) 在 \(AC\) 上，\(BF\perp AO\).



\begin{enumerate}
\item 证明：\(EF // \text{平面 }ADO\)；
\item 证明：\(\text{平面 }ADO\perp\text{平面 }BEF\)；
\item 求二面角 \(D-AO-C\) 的正弦值.
\end{enumerate}

\bitmapfigure[max width=6.0cm]{img/2023/national_paper_b_science/q19_fig1.png}

\end{problem}

\begin{solution}
（1）连接 \(DE\)，\(OF\)，设 \(AF=tAC\)，则 \(\overrightarrow{BF}=\overrightarrow{BA}+\overrightarrow{AF}=(1-t)\overrightarrow{BA}+t\overrightarrow{BC}\)，\(\overrightarrow{AO}=-\overrightarrow{BA}+\frac12\overrightarrow{BC}\). 因为 \(BF\perp AO\)，所以
\[
\overrightarrow{BF}\cdot\overrightarrow{AO}=\left[ (1-t)\overrightarrow{BA}+t\overrightarrow{BC} \right]\cdot\left( -\overrightarrow{BA}+\frac12\overrightarrow{BC} \right)=4(t-1)+4t=0
\]
解得 \(t=\frac12\)，则 \(F\) 为 \(AC\) 的中点.
由 \(D\)，\(E\)，\(O\)，\(F\) 分别为 \(PB\)，\(PA\)，\(BC\)，\(AC\) 的中点，\(DE//AB\)，\(DE=\frac12AB\)，\(OF//AB\)，\(OF=\frac12AB\)，即 \(DE//OF\)，且 \(DE=OF\)，所以四边形 \(ODEF\) 为平行四边形，从而 \(EF//DO\)，\(EF=DO\). 又 \(EF\not\subset\text{平面 }ADO\)，\(DO\subset\text{平面 }ADO\)，所以 \(EF // \text{平面 }ADO\).

（2）由（1）可知 \(EF//OD\)，则 \(AO=\sqrt6\)，\(DO=\frac{\sqrt6}{2}\)，\(AD=\sqrt5 DO=\frac{\sqrt{30}}{2}\). 因此
\[
OD^2+AO^2=\frac{6}{4}+6=\frac{15}{2}=AD^2
\]
故 \(OD\perp AO\)，于是 \(EF\perp AO\).
又 \(AO\perp BF\)，且 \(BF\cap EF=F\)，\(BF\)，\(EF\subset\text{平面 }BEF\)，所以 \(AO\perp\text{平面 }BEF\). 又 \(AO\subset\text{平面 }ADO\)，故 \(\text{平面 }ADO\perp\text{平面 }BEF\).

（3）过点 \(O\) 作 \(OH//BF\) 交 \(AC\) 于点 \(H\)，设 \(AD\cap BE=G\).
由 \(AO\perp BF\)，得 \(HO\perp AO\)，且 \(FH=\frac13AH\). 又由（2）知 \(OD\perp AO\)，则 \(\angle DOH\) 为二面角 \(D-AO-C\) 的平面角.
因为 \(D\)，\(E\) 分别为 \(PB\)，\(PA\) 的中点，所以 \(G\) 为 \(\triangle PAB\) 的重心，从而 \(DG=\frac13AD\)，\(GE=\frac13BE\). 又 \(FH=\frac13AH\)，故 \(DH=\frac32GF\).
由
\[
\cos\angle ABD=\frac{AB^2+BD^2-AD^2}{2AB\cdot BD}=\frac{4+\frac64-\frac{30}{4}}{2\times2\times\frac{\sqrt6}{2}}=-\frac{1}{\sqrt6}
\]
又因为 \(B\)，\(D\)，\(P\) 三点共线，所以 \(\angle ABD=\angle ABP\)，故
\[
\cos\angle ABP=\frac{AB^2+BP^2-PA^2}{2AB\cdot BP}=\frac{4+6-PA^2}{2\times2\times\sqrt6}=-\frac{1}{\sqrt6}
\]
解得 \(PA=\sqrt{14}\). 又有 \(\cos\angle BAP=\frac{4+14-6}{2\times2\times\sqrt{14}}=\frac{3}{\sqrt{14}}\).
因为 \(E\) 为 \(PA\) 的中点，所以 \(AE=\frac12PA=\frac{\sqrt{14}}{2}\)，\(\angle BAE=\angle BAP\). 在 \(\triangle ABE\) 中，
\[
BE^2=AB^2+AE^2-2AB\cdot AE\cos\angle BAE=4+\frac{14}{4}-2\times2\times\frac{\sqrt{14}}{2}\times\frac{3}{\sqrt{14}}=\frac32
\]
所以 \(BE=\frac{\sqrt6}{2}\).
由 \(F\) 为 \(AC\) 的中点，且 \(\triangle ABC\) 为直角三角形，可得
\[
BF=\frac12AC=\frac12\sqrt{AB^2+BC^2}=\sqrt3
\]
又由（1）知 \(EF=DO=\frac{\sqrt6}{2}\)，因此
\[
BE^2+EF^2=\frac32+\frac32=3=BF^2
\]
所以 \(BE\perp EF\). 又因为 \(G\) 在 \(BE\) 上，所以 \(GE\perp EF\)，于是 \(GF^2=\left( \frac13\times\frac{\sqrt6}{2} \right)^2+\left( \frac{\sqrt6}{2} \right)^2=\frac53\).
故 \(GF=\frac{\sqrt{15}}{3}\)，\(DH=\frac32\times\frac{\sqrt{15}}{3}=\frac{\sqrt{15}}{2}\).
在 \(\triangle DOH\) 中，\(OH=\frac12BF=\frac{\sqrt3}{2}\)，\(OD=\frac{\sqrt6}{2}\)，\(DH=\frac{\sqrt{15}}{2}\)，所以
\(\cos\angle DOH=\frac{\frac64+\frac34-\frac{15}{4}}{2\times\frac{\sqrt6}{2}\times\frac{\sqrt3}{2}}=-\frac{\sqrt2}{2}\).
\(\sin\angle DOH=\sqrt{1-\left( -\frac{\sqrt2}{2} \right)^2}=\frac{\sqrt2}{2}\). 所以二面角 \(D-AO-C\) 的正弦值为 \(\frac{\sqrt2}{2}\).
\end{solution}

\begin{problem}
已知椭圆
\(\frac{y^2}{a^2}+\frac{x^2}{b^2}=1\)（\(a>b>0\)），其离心率是 \(\frac{\sqrt5}{3}\)，点 \(A(-2,0)\) 在 \(C\) 上.

\begin{enumerate}
\item 求 \(C\) 的方程；
\item 过点 \((-2,3)\) 的直线交 \(C\) 于 \(P\)，\(Q\) 两点，直线 \(AP\)，\(AQ\) 与 \(y\) 轴的交点分别为 \(M\)，\(N\)，证明：线段 \(MN\) 的中点为定点.
\end{enumerate}

\end{problem}

\begin{solution}
（1）由题意可得
\(a^2=b^2+c^2\)，\(\dfrac{c}{a}=\dfrac{\sqrt5}{3}\)，
解得
\(a=3\)，\(b=2\)，\(c=\sqrt5\).
所以椭圆 \(C\) 的方程为
\(\frac{y^2}{9}+\frac{x^2}{4}=1\).

（2）由题意可知：直线 \(PQ\) 的斜率存在，设
\(PQ:\ y=k(x+2)+3\)，\(P(x_1,y_1)\)，\(Q(x_2,y_2)\).
联立方程
\(y=k(x+2)+3\)，\(\dfrac{y^2}{9}+\dfrac{x^2}{4}=1\)，
消去 \(y\) 得
\((4k^2+9)x^2+8k(2k+3)x+16(k^2+3k)=0\).
则
\[
\Delta=64k^2(2k+3)^2-64(4k^2+9)(k^2+3k)=-1728k>0
\]
解得 \(k<0\)，并且
\(x_1+x_2=-\frac{8k(2k+3)}{4k^2+9}\)，\(x_1x_2=\frac{16(k^2+3k)}{4k^2+9}\).
因为 \(A(-2,0)\)，则直线 \(AP\) 的方程为
\(y=\frac{y_1}{x_1+2}(x+2)\).
令 \(x=0\)，解得
\(y=\frac{2y_1}{x_1+2}\)，
即
\(M\left( 0,\frac{2y_1}{x_1+2} \right)\)，
则直线 \(AQ\) 的方程为
\(y=\frac{y_2}{x_2+2}(x+2)\).
令 \(x=0\)，解得
\(y=\frac{2y_2}{x_2+2}\)，
即
\(N\left( 0,\frac{2y_2}{x_2+2} \right)\).
于是 \(\frac{\frac{2y_1}{x_1+2}+\frac{2y_2}{x_2+2}}{2}=\frac{[k(x_1+2)+3](x_2+2)+[k(x_2+2)+3](x_1+2)}{(x_1+2)(x_2+2)}\).
化简得
\(\frac{2kx_1x_2+(4k+3)(x_1+x_2)+4(2k+3)}{x_1x_2+2(x_1+x_2)+4}=3\).
所以线段 \(MN\) 的中点是定点 \((0,3)\).
\end{solution}

\begin{problem}
已知
\(f(x)=\left( \frac1x+a \right)\ln(1+x)\).

\begin{enumerate}
\item 当 \(a=-1\) 时，求曲线 \(y=f(x)\) 在点 \((1,f(1))\) 处的切线方程；
\item 是否存在 \(a\)，\(b\)，使得曲线 \(y=f\left( \frac1x \right)\) 关于直线 \(x=b\) 对称，若存在，求 \(a\)，\(b\) 的值，若不存在，说明理由；
\item 若 \(f(x)\) 在 \((0,+\infty)\) 存在极值，求 \(a\) 的取值范围.
\end{enumerate}

\end{problem}

\begin{solution}
（1）当 \(a=-1\) 时，
\(f(x)=\left( \frac1x-1 \right)\ln(x+1)\)，
则
\(f'(x)=-\frac{1}{x^2}\ln(x+1)+\left( \frac1x-1 \right)\frac{1}{x+1}\).
据此可得 \(f(1)=0\)，\(f'(1)=-\ln2\)，
所以曲线在 \((1,f(1))\) 处的切线方程为
\(y-0=-\ln2(x-1)\)，
即
\((\ln2)x+y-\ln2=0\).

（2）由函数的解析式可得
\(f\left( \frac1x \right)=(x+a)\ln\left( \frac1x+1 \right)\).
函数的定义域满足
\(\frac1x+1=\frac{x+1}{x}>0\)，
即定义域为 \((-\infty,-1)\cup(0,+\infty)\)，它关于直线 \(x=-\frac12\) 对称，
故由题意可得
\(b=-\frac12\).
由对称性可知
\(f\left( -\frac12+m \right)=f\left( -\frac12-m \right)\)（\(m>\frac12\)）.
取 \(m=\frac32\)，可得 \(f(1)=f(-2)\)，
即
\((a+1)\ln2=(a-2)\ln\frac12\)，
所以
\(a+1=2-a\)，
解得
\(a=\frac12\).
经检验 \(a=\frac12\)，\(b=-\frac12\) 满足题意.

（3）由函数的解析式可得
\(f'(x)=-\frac{1}{x^2}\ln(x+1)+\left( \frac1x+a \right)\frac{1}{x+1}\).
由 \(f(x)\) 在区间 \((0,+\infty)\) 存在极值点，则 \(f'(x)\) 在区间 \((0,+\infty)\) 上存在变号零点.
令
\(-\frac{1}{x^2}\ln(x+1)+\left( \frac1x+a \right)\frac{1}{x+1}=0\)，
则
\(-(x+1)\ln(x+1)+(x+ax^2)=0\).
记
\(g(x)=ax^2+x-(x+1)\ln(x+1)\)，
则 \(f(x)\) 在区间 \((0,+\infty)\) 存在极值点，等价于 \(g(x)\) 在区间 \((0,+\infty)\) 上存在变号零点.
有
\(g'(x)=2ax-\ln(x+1)\)，\(g''(x)=2a-\frac{1}{x+1}\).

当 \(a\le0\) 时，\(g'(x)<0\)，\(g(x)\) 在区间 \((0,+\infty)\) 上单调递减，
此时 \(g(x)<g(0)=0\)，\(g(x)\) 在区间 \((0,+\infty)\) 上无零点，不合题意.

当 \(a\ge\frac12\) 时，因为 \(\frac{1}{x+1}<1\)，所以 \(g''(x)>0\)，\(g'(x)\) 在区间 \((0,+\infty)\) 上单调递增，
故 \(g'(x)>g'(0)=0\)，\(g(x)\) 在区间 \((0,+\infty)\) 上单调递增，
从而 \(g(x)>g(0)=0\)，\(g(x)\) 在区间 \((0,+\infty)\) 上无零点，不符合题意.

当 \(0<a<\frac12\) 时，由
\(g''(x)=2a-\frac{1}{x+1}=0\)
可得
\(x=\frac{1}{2a}-1\).
当 \(x\in\left( 0,\frac{1}{2a}-1 \right)\) 时，\(g''(x)<0\)，\(g'(x)\) 单调递减；
当 \(x\in\left( \frac{1}{2a}-1,+\infty \right)\) 时，\(g''(x)>0\)，\(g'(x)\) 单调递增，
故 \(g'(x)\) 的最小值为
\(g'\left( \frac{1}{2a}-1 \right)=1-2a+\ln2a\).
令 \(m(x)=1-x+\ln x\)（\(0<x<1\)），则
\(m'(x)=\frac{-x+1}{x}>0\)，
所以 \(m(x)\) 在定义域内单调递增，且 \(m(x)<m(1)=0\)，
据此可得 \(1-x+\ln x<0\) 恒成立，
于是
\(g'\left( \frac{1}{2a}-1 \right)=1-2a+\ln2a<0\).
令 \(h(x)=\ln x-x^2+x\)（\(x>0\)），则
\(h'(x)=\frac{-2x^2+x+1}{x}\).
当 \(x\in(0,1)\) 时，\(h'(x)>0\)，\(h(x)\) 单调递增；
当 \(x\in(1,+\infty)\) 时，\(h'(x)<0\)，\(h(x)\) 单调递减.
故 \(h(x)\le h(1)=0\)，即
\(\ln x\le x^2-x\)
（取等条件为 \(x=1\)）.
所以
\(g'(x)=2ax-\ln(x+1)>2ax-\left[ (x+1)^2-(x+1) \right]=2ax-(x^2+x)\)，
从而
\(g'(2a-1)>2a(2a-1)-\left[ (2a-1)^2+(2a-1) \right]=0\).
又注意到 \(g'(0)=0\)，根据零点存在性定理可知：\(g'(x)\) 在区间 \((0,+\infty)\) 上存在唯一零点 \(x_0\).
当 \(x\in(0,x_0)\) 时，\(g'(x)<0\)，\(g(x)\) 单调减；
当 \(x\in(x_0,+\infty)\) 时，\(g'(x)>0\)，\(g(x)\) 单调递增，
所以
\(g(x_0)<g(0)=0\).
令
\(n(x)=\ln x-\frac12\left( x-\frac1x \right)\)，
则
\(n'(x)=\frac1x-\frac12\left( 1+\frac{1}{x^2} \right)=\frac{-(x-1)^2}{2x^2}\le0\).
故 \(n(x)\) 单调递减，且 \(n(1)=0\)，
所以当 \(x\in(1,+\infty)\) 时，
\(\ln x-\frac12\left( x-\frac1x \right)<0\)，
从而
\(\ln x<\frac12\left( x-\frac1x \right)\).
于是 \(g(x)=ax^2+x-(x+1)\ln(x+1)>ax^2+x-(x+1)\times\frac12\left[ (x+1)-\frac{1}{x+1} \right]=\left( a-\frac12 \right)x^2+\frac12\).
令
\(\left( a-\frac12 \right)x^2+\frac12=0\)
得
\(x_2=\sqrt{\frac{1}{1-2a}}\)，
故
\(g\left( \sqrt{\frac{1}{1-2a}} \right)>0\).
所以函数 \(g(x)\) 在区间 \((0,+\infty)\) 上存在变号零点，符合题意.
综合上面可知：实数 \(a\) 的取值范围是 \(\left( 0,\frac12 \right)\).
\end{solution}

\begin{problem}
在直角坐标系 \(xOy\) 中，以坐标原点 \(O\) 为极点，\(x\) 轴正半轴为极轴建立极坐标系，曲线 \(C_1\) 的极坐标方程为 \(\rho=2\sin\theta\)（\(\frac{\pi}{4}\le\theta\le\frac{\pi}{2}\)），曲线 \(C_2\) 的参数方程为 \(x=2\cos\alpha\)，\(y=2\sin\alpha\)（\(\alpha\) 为参数，\(\frac{\pi}{2}<\alpha<\pi\)）.

\begin{enumerate}
\item 写出 \(C_1\) 的直角坐标方程；
\item 若直线 \(y=x+m\) 既与 \(C_1\) 没有公共点，也与 \(C_2\) 没有公共点，求 \(m\) 的取值范围.
\end{enumerate}

\end{problem}

\begin{solution}
（1）因为 \(\rho=2\sin\theta\)，即 \(\rho^2=2\rho\sin\theta\)，可得 \(x^2+y^2=2y\). 整理得 \(x^2+\left( y-1 \right)^2=1\).
表示以 \((0,1)\) 为圆心，半径为 1 的圆.
又因为 \(x=\rho\cos\theta=2\sin\theta\cos\theta=\sin2\theta\)，\(y=\rho\sin\theta=2\sin^2\theta=1-\cos2\theta\)，且 \(\frac{\pi}{4}\le\theta\le\frac{\pi}{2}\)，则 \(\frac{\pi}{2}\le2\theta\le\pi\)，所以 \(x=\sin2\theta\in[0,1]\)，\(y=1-\cos2\theta\in[1,2]\).
故
\(C_1:\ x^2+\left( y-1 \right)^2=1\)，\(x\in[0,1]\)，\(y\in[1,2]\)

（2）因为 \(C_2\) 的参数方程为 \(x=2\cos\alpha\)，\(y=2\sin\alpha\)（\(\frac{\pi}{2}<\alpha<\pi\)），整理得 \(x^2+y^2=4\).
表示圆心为 \(O(0,0)\)，半径为 2，且位于第二象限的圆弧.
如图所示，若直线 \(y=x+m\) 过 \((1,1)\)，则 \(1=1+m\)，解得 \(m=0\).
若直线 \(y=x+m\)，即 \(x-y+m=0\) 与 \(C_2\) 相切，则 \(\frac{|m|}{\sqrt2}=2\)，又 \(m>0\)，解得 \(m=2\sqrt2\).
若直线 \(y=x+m\) 与 \(C_1\)，\(C_2\) 均没有公共点，则 \(m>2\sqrt2\) 或 \(m<0\)，即实数 \(m\) 的取值范围为 \((-\infty,0)\cup(2\sqrt2,+\infty)\).
\end{solution}

\begin{problem}
已知 \(f(x)=2|x|+|x-2|\).

\begin{enumerate}
\item 求不等式 \(f(x)\le6-x\) 的解集；
\item 在直角坐标系 \(xOy\) 中，求由不等式 \(f(x)\le y\) 与 \(x+y-6\le0\) 所确定的平面区域的面积.
\end{enumerate}

\end{problem}

\begin{solution}
（1）依题意，
\[
f(x)=\begin{cases}
3x-2,&x>2,\\
x+2,&0\le x\le2,\\
-3x+2,&x<0.
\end{cases}
\]
不等式 \(f(x)\le6-x\) 化为
\[
\begin{cases}
x>2,\\
3x-2\le6-x
\end{cases}
\quad\text{或}\quad
\begin{cases}
0\le x\le2,\\
x+2\le6-x
\end{cases}
\quad\text{或}\quad
\begin{cases}
x<0,\\
-3x+2\le6-x.
\end{cases}
\]
解得第一组无解，第二组得 \(0\le x\le2\)，第三组得 \(-2\le x<0\). 因此原不等式的解集为 \([-2,2]\).

（2）作出不等式组 \(f(x)\le y\) 与 \(x+y-6\le0\) 所表示的平面区域，如图中阴影 \(\triangle ABC\).
由 \(y=-3x+2\)，\(x+y=6\)，
解得 \(A(-2,8)\)；由
\(y=x+2\)，\(x+y=6\)，
解得 \(C(2,4)\)；又 \(B(0,2)\)，\(D(0,6)\).
所以 \(S_{\triangle ABC}=\frac12|BD|\times|x_C-x_A|=\frac12|6-2|\times|2-(-2)|=8\).
\end{solution}
